%------------------------------------------------------------------------------%
\documentclass[a4paper,12pt,fleqn]{article}

\usepackage{gaal}

\ShowAnswers

%------------------------------------------------------------------------------%
\begin{document}

\makeHeader{GAAL}{Prova 2 -- Turma 4}{23/09/2026}

% 01
%------------------------------------------------------------------------------%
\exercise{20\%}
Calcule determinante de
\(
  A =
  \begin{pmatrix}
    1 & 2 & 1 & 0 \\
    2 & 1 & 0 & 1 \\
    1 & 0 & 2 & 1 \\
    0 & 1 & 1 & 2
  \end{pmatrix}
\)
e determine se a matriz é invertível

\begin{answer}
\begin{align*}
  \det(A)
  & = a_{11}(-1)^{1+1} M_{11}
    + a_{12}(-1)^{1+2} M_{12}
    + a_{13}(-1)^{1+3} M_{13}
    + a_{14}(-1)^{1+4} M_{14} \\
  & = 1 M_{11} - 2 M_{12} + 1 M_{13} - 0 M_{14}   \\
  & = M_{11} - 2 M_{12} + M_{13}
\end{align*}
\[
  M_{11}
  = \begin{vmatrix}
      1 & 0 & 1 \\
      0 & 2 & 1 \\
      1 & 1 & 2
    \end{vmatrix}
  = 1 (-1)^{1+1}
    \begin{vmatrix}
      2 & 1 \\
      1 & 2
    \end{vmatrix}
    + 0
    + 1 (-1)^{1+3}
    \begin{vmatrix}
      0 & 2 \\
      1 & 1
    \end{vmatrix}
  = (4-1)+(0-2)
  = 3 - 2
  = 1
\]
\[
  M_{12}
  = \begin{vmatrix}
    2 & 0 & 1 \\
    1 & 2 & 1 \\
    0 & 1 & 2
  \end{vmatrix}
  = 2  (-1)^{1+1}
  \begin{vmatrix}
    2 & 1 \\
    1 & 2
  \end{vmatrix}
  + 0
  + 1 (-1)^{1+2}
  \begin{vmatrix}
    1 & 2 \\
    0 & 1
  \end{vmatrix}
  = 2(4-1)+(1-0)
  = 6 + 1
  = 7
\]
\[
  M_{13}
  = \begin{vmatrix}
    2 & 1 & 1 \\
    1 & 0 & 1 \\
    0 & 1 & 2
  \end{vmatrix}
  = 1 (-1)^{2+1}
  \begin{vmatrix}
    1 & 1 \\
    1 & 2
  \end{vmatrix}
  + 0
  - 1 (-1)^{2+3}
  \begin{vmatrix}
    2 & 1 \\
    0 & 1
  \end{vmatrix}
  = -(4-0) - (2-0)
  = -4 - 2
  = -6
\]
\[
  \det(A) = 1 -2 \times 7 -6 = -19
\]

Como
\(
  \det(A) \neq 0,
\)
$A$ é invertível
\end{answer}

% 02
%------------------------------------------------------------------------------%
\exercise{25\%}
Calcule a projeção ortogonal de
\(
  \vec{r} = \vetortri{1}{2}{-1}
\)
em
\(
  \vec{s} = \vetortri{2}{-1}{3}
\)

\begin{answer}
  \[
    \vec{r}\cdot\vec{s}
    = \vetortri{1}{2}{-1} \cdot \vetortri{2}{-1}{3}
    = 1 \times 2 + 2(-1) + (-1)3
    = 2 - 2 - 3
    = -3
  \]
  \[
    \vec{s}\cdot\vec{s}
    = \vetortri{2}{-1}{3} \cdot \vetortri{2}{-1}{3}
    = 2^2 + (-1)^2 + 3^2
    = 4 + 1 + 9
    = 14
  \]
  \[
    \projuv{r}{s}
    = \frac{\vec{r}\cdot\vec{s}}{\vec{s}\cdot\vec{s}} \, \vec{s}
    = \frac{-3}{14} \vetortri{2}{-1}{3}
    = \Vetortri{\frac{-3}{7}}{\frac{3}{14}}{\frac{-9}{14}}
  \]
\end{answer}

% 03
%------------------------------------------------------------------------------%
\exercise{20\%}
Encontre o vetor unitário na direção e sentido de
\(
  \vec{\psi} = 2 \vetortri{2}{-1}{3} - 3 \vetortri{1}{2}{-1}
\)

\begin{answer}
  \[
    \vec{\psi}
    = 2\vetortri{2}{-1}{3} - 3\vetortri{1}{2}{-1}
    = \vetortri{4}{-2}{6} + \vetortri{-3}{-6}{3}
    = \vetortri{1}{-8}{9}
  \]
  \[
    \norm{\vec{\psi}}
    = \norm{\vetortri{1}{-8}{9}}
    = \sqrt{1^2+(-8)^2+9^2}
    = \sqrt{1+64+81}
    = \sqrt{146}
  \]
  \[
    \hat{\psi}
    = \frac{\vec{\psi}}{\norm{\vec{\psi}}}
    = \frac{1}{\sqrt{146}} \vetortri{1}{-8}{9}
    = \Vetortri{\frac{1}{\sqrt{146}}}{\frac{-8}{\sqrt{146}}}{\frac{9}{\sqrt{146}}}
  \]
\end{answer}

% 04
%------------------------------------------------------------------------------%
\exercise{20\%}
Calcule
\(
  \norm{(1,  2, 1) \times (2, -1, 3)}
\)
\begin{answer}
  \begin{align*}
    \vec{\psi}
    & = (1,  2, 1) \times (2, -1, 3) \\
    & = \begin{vmatrix}
          \vec{i} & \vec{j} & \vec{k} \\
          1        & 2        & 1       \\
          2        & -1       & 3
        \end{vmatrix}
    \\
    & = \vec{i}\, (-1)^{1+1}
        \begin{vmatrix}
          2 & 1 \\
          -1 & 3
        \end{vmatrix}
      + \vec{j}\, (-1)^{1+2}
        \begin{vmatrix}
          1 & 1 \\
          2 & 3
        \end{vmatrix}
      + \vec{k}\, (-1)^{1+3}
        \begin{vmatrix}
          1 & 2 \\
          2 & -1
        \end{vmatrix}
    \\
    & = \vec{i} (6 + 1)
      - \vec{j} (3 - 2)
      + \vec{k} (-1 -4) \\
    & = 7\vec{i} - \vec{j} -5\vec{k} \\
    & = \vetortri{7}{-1}{-5}
  \end{align*}
  \[
    \norm{\vec{\psi}}
    = \norm{\vetortri{7}{-1}{5}}
    = \sqrt{7^2 + (-1)^2 + 5^2}
    = \sqrt{49 + 1 + 25}
    = \sqrt{75}
  \]
\end{answer}

% 05
%------------------------------------------------------------------------------%
\exercise{20\%}
calcule o produto misto dos vetores
\(
  \vec{p} = \vetortri{1}{2}{0}
\),
\(
  \vec{q} = \vetortri{2}{-1}{1}
\) e
\(
  \vec{r} = \vetortri{1}{0}{3}
\)

\begin{answer}
\begin{align*}
  \left(\vec{p} \times \vec{q}\right) \cdot \vec{r}
  & = \begin{vmatrix}
    1 & 2  & 0 \\
    2 & -1 & 1 \\
    1 & 0  & 3
  \end{vmatrix}
  \\
  & = 1
    \begin{vmatrix}
      -1 & 1 \\
      0  & 3
    \end{vmatrix}
    - 2
    \begin{vmatrix}
      2 & 1 \\
      1 & 3
    \end{vmatrix}
    + 0
  \\
  & = 1(-3-0) -2(6-1) \\
  & = -3 - 10 \\
  & = -13
\end{align*}
\end{answer}

%------------------------------------------------------------------------------%
\end{document}
%------------------------------------------------------------------------------%

